\documentclass[a4paper,11pt,final]{article}
% Pour une impression recto verso, utilisez plutôt ce documentclass :
%\documentclass[a4paper,11pt,twoside,final]{article}

\usepackage[english,francais]{babel}
\usepackage[utf8]{inputenc}
\usepackage[T1]{fontenc}
\usepackage[pdftex]{graphicx}
\usepackage{setspace}
\usepackage{hyperref}
\usepackage[french]{varioref}
\usepackage{geometry}
\geometry{hmargin=2.5cm,vmargin=2.5cm}

\newcommand{\reporttitle}{Rapport de Stage}     % Titre
\newcommand{\reportauthor}{Thibaud \textsc{Courtoison}} % Auteur
\newcommand{\reportsubject}{Stage de fin d'étude} % Sujet
\newcommand{\HRule}{\rule{\linewidth}{0.5mm}}
\setlength{\parskip}{1ex} % Espace entre les paragraphes

\hypersetup{
    pdftitle={\reporttitle},%
    pdfauthor={\reportauthor},%
    pdfsubject={\reportsubject},%
    pdfkeywords={rapport} {vos} {mots} {clés}
}

\begin{document}
  % Inspiré de http://en.wikibooks.org/wiki/LaTeX/Title_Creation

\begin{titlepage}

\begin{center}

% Images entreprise & association
\begin{minipage}[t]{0.48\textwidth}
  \begin{flushleft}
    \includegraphics [height=30mm]{logo-libertic.png} \\[0.2cm]
    \begin{spacing}{1.5}
      \textsc{\LARGE Association\\LiberTIC}
    \end{spacing}
  \end{flushleft}
\end{minipage}
\begin{minipage}[t]{0.48\textwidth}
  \begin{flushright}
    \includegraphics [height=30mm]{logo-polypodes.png} \\[0.5cm]
    \textsc{\LARGE Entreprise\\Les Polypodes}
  \end{flushright}
\end{minipage} %\\[0.8cm]

\vfill

% Titre Rapport
\textsc{\Large \reportsubject}\\[0.5cm]
\HRule \\[0.6cm]
{\huge \bfseries \reporttitle}\\[0.4cm]
\HRule \\[1.0cm]

% Auteur et encadrant
\begin{minipage}[t]{0.3\textwidth}
  \begin{flushleft} \large
    \emph{Auteur :}\\
    \reportauthor
  \end{flushleft}
\end{minipage}
\begin{minipage}[t]{0.6\textwidth}
  \begin{flushright} \large
    \emph{Responsables :} \\
    M.~Ronan \textsc{Guilloux} \\
    M.~Nicolas \textsc{Hernandez}
  \end{flushright}
\end{minipage}


\vfill

\begin{minipage}[t]{0.48\textwidth}

  \begin{flushleft}
    \includegraphics [height=30mm]{logo-univ.jpg} \\[0.2cm]
    \begin{spacing}{1.5}
      \textsc{\LARGE Université de Nantes}
    \end{spacing}
  \end{flushleft}

\end{minipage}
\begin{minipage}[t]{0.48\textwidth}

  \begin{flushright}
    \includegraphics [height=30mm]{logo-iut.jpg} \\[0.2cm]
    \begin{spacing}{1.5}
      \textsc{\LARGE IUT de Nantes}
    \end{spacing}
  \end{flushright}

\end{minipage}
%\vfill

{\large Du 13 avril 2015 au 19 juin 2015}

\end{center}

\end{titlepage}

  \cleardoublepage % Dans le cas du recto verso, ajoute une page blanche si besoin
  \tableofcontents % Table des matières
  \sloppy          % Justification moins stricte : des mots ne dépasseront pas des paragraphes
  \cleardoublepage
  \section*{Remerciements}
\addcontentsline{toc}{section}{Remerciements}

Je tiens tout d'abord à remercier les personnes qui m'ont accompagnés durant ces deux années de DUT.\\

\textit{M. Ronan Guilloux} et toute l'équipe des Polypodes, pour m'avoir accordé leurs confiance en m'acceptant dans leur équipe et pour m'avoir accompagné durant ces deux mois et demi. Leur patience et leur aide ont fait de ce stage une expérience particulièrement constructive.\\

Et bien sûr, toute l'équipe pédagogique du département informatique de l'IUT de Nantes, qui m'ont fait découvrir les mondes passionnants de l'informatique, de la gestion, du droit et de la communication. Si ces deux années peuvent s'assimiler à une réussite, c'est en très grande partie grâce à eux.\\

\nb{Refact this}

  \cleardoublepage
  \section*{Introduction} % Pas de numérotation
\addcontentsline{toc}{section}{Introduction} % Ajout dans la table des matières

Ce document est un exemple de rapport écrit en \LaTeX.

Écrit par Thibaud Courtoison grâce à \url{http://blog.hikoweb.net/}.

Pour rajouter une correction ou une note, utiliser la commande \\ronan. Example:

\textbackslash{}ronan\{Tu as oublié de parler de XXX\}

va afficher:

\ronan{Tu as oublié de parler de XXX}
  \cleardoublepage
  \section{La première section}

  \cleardoublepage
  \section{La deuxième section}

  \cleardoublepage
  \section*{Conclusion}
\addcontentsline{toc}{section}{Conclusion}

Pour conclure, avec \LaTeX{} on obtient un rendu impeccable mais il faut s'investir pour le prendre en main.

  \cleardoublepage
  \phantomsection\addcontentsline{toc}{section}{Références}

\renewcommand{\refname}{Webographie}
\newcounter{firstbib}

\begin{thebibliography}{ABC}

    \bibitem{cartopendata} LiberTIC. \emph{OpenData Map}. URL: \url{http://www.opendata-map.org/}. Consulté le 1 juin 2015.

    \bibitem{gitimmersion} Gitimmersion. \emph{Best GIT tuto ever -50 steps}. URL: \url{http://gitimmersion.com/}. Consulté le 13 avril 2015.

    \bibitem{successfulgit} Vincent Driessen. \emph{A successful Git branching model}. Janvier 2010. URL: \url{http://nvie.com/posts/a-successful-git-branching-model/}. Consulté le 13 avril 2015.

    \bibitem{bestpracticephp} Brian Fenton. \emph{Best Practices for Modern PHP Development}. URL: \url{https://www.airpair.com/php/posts/best-practices-for-modern-php-development}. Consulté le 14 avril 2015.

    \bibitem{stupidvssolid} Hugo Hamon. \emph{Votre code est STUPID ? Rendez le SOLID}. Décembre 2013. URL: \url{http://afsy.fr/avent/2013/02-principes-stupid-solid-poo}. Consulté le 14 avril 2015.

    \bibitem{symfonybook} Sensio Labs. \emph{The Symfony Book}. URL: \url{http://symfony.com/doc/current/book/index.html}. Consulté plusieurs fois du 15 avril 2015 au 19 juin 2015.

    \bibitem{restsymfony} William Durand. \emph{REST APIs with Symfony2: The Right Way}. Août 2012. URL: \url{http://williamdurand.fr/2012/08/02/rest-apis-with-symfony2-the-right-way/}. Consulté le 22 mai 2015.

    \setcounter{firstbib}{\value{enumiv}}
\end{thebibliography}

\renewcommand{\refname}{RFCs}

\begin{thebibliography}{ABC}
    \setcounter{enumiv}{\value{firstbib}}

    \bibitem{rfc4918} RFC 4918. \emph{HTTP Extensions for Web Distributed Authoring and Versioning}. URL: \url{http://tools.ietf.org/html/rfc4918}.

    \bibitem{rfc4791} RFC 4791. \emph{Calendaring Extensions to WebDAV}. URL: \url{http://tools.ietf.org/html/rfc4791}.

    \bibitem{rfc6352} RFC 6352. \emph{vCard Extensions to WebDAV}. URL: \url{http://tools.ietf.org/html/rfc6352}.

    \bibitem{rfc5545} RFC 5545. \emph{Internet Calendaring and Scheduling Core Object Specification}. URL: \url{http://tools.ietf.org/html/rfc5545}

    \bibitem{rfc6350} RFC 6350. \emph{vCard Format Specification}. URL: \url{http://tools.ietf.org/html/rfc6350}.

\end{thebibliography}

\end{document}